% The labbook: the logical steps of this project, for a reader — a colleague,
% a collaborator, yourself in six months. Written by the agent in the
% researcher's own voice (labbook/voice.md), one \section per step.
%
% The boxes, figures and tables come from the claim ledger through
% claims.gen.tex (`labbook.py gen`). Never copy a claim's text in here:
% \claimbox{<id>} shows it exactly as the ledger has it.
\documentclass[11pt,a4paper]{article}
\usepackage[english]{babel}
\usepackage{berry-labbook}
\input{claims.gen.tex}

\title{@TITLE@}
\author{@AUTHOR@}
\date{\today}

\begin{document}
\maketitle

\section*{Where things stand}
% Three or four sentences: the question, what has been established so far,
% and what is open. Rewritten as the project moves.

\berryoverview

% --- the steps ---------------------------------------------------------------
%
% \section{<what this step found, in a few words>}
% Why this step: the question it answered, and why it came next.
% What was done: the method, in plain words; what was held fixed.
% \claimbox{<claim id>}
% \claimfigures{<claim id>}
% What it means for the next step. Papers checked: \sourcebox{<sources/…json>}

% --- end -----------------------------------------------------------------------

\section*{Open points and follow-up questions}
\begin{itemize}
  \item % What is not settled, what to try next, what to ask.
\end{itemize}

% \section*{Notes}
% Only what the researcher asked to have written down beyond the claims:
% context, decisions, remarks. Never in place of a claim. Remove the % to use it.

\appendix
\section{Conventions}
\berryconventions

\end{document}
